% ============================================================
% Introduction § 1 + § 2 (mechanism) — DRAFT for coworker review.
%
% This file is BOTH a standalone-compilable preview AND a fragment
% you can \input{} into the main manuscript. To use as a fragment,
% delete the documentclass / begin{document} / end{document} block.
%
% Length budgets:
%   § 1 (the hook):       ~250 words
%   § 2 (the mechanism):  ~400 words
%
% Tone: methodological-correction first, model second.
% Editorial guidance: this is what an editor at NComms reads
%   in the first 90 seconds. Every sentence has to do work.
%
% Markers:
%   [TODO/CW]  = needs coworker confirmation or supplied value
%   [PHASE B]  = number from the still-running audit matrix
% ============================================================
\documentclass[11pt]{article}
\usepackage[a4paper,margin=2.5cm]{geometry}
\usepackage{amsmath}\usepackage{amssymb}\usepackage{microtype}
\usepackage{hyperref}
\setlength{\parskip}{6pt}
\title{[Draft fragment] Test-edge leakage inflates seven years\\of signed-link prediction benchmarks; a strict protocol\\and audit restore comparability}
\author{[Authors]}
\date{}
\begin{document}\maketitle

% =========================================================
% § 1 — The editorial hook (~250 words)
% =========================================================
\section{Introduction}

Signed networks --- where edges carry a positive (cooperation, trust,
agreement) or negative (interference, distrust, disagreement) sign
in addition to topological information --- have become the canonical
benchmark domain for evaluating sign-aware graph neural networks.
Five publicly available datasets (Bitcoin-Alpha, Bitcoin-OTC,
Slashdot, Epinions, Reddit Hyperlinks) have anchored a seven-year
methodological trajectory: SGCN~\cite{sgcn2018}, SiGAT~\cite{sigat2019},
SGCL~\cite{sgcl2022}, SiGformer~\cite{sigformer2023},
SE-SGformer~\cite{sesgformer2024}, and DADSGNN~\cite{dadsgnn2025},
with peer-reviewed AUROC ceilings rising from $\sim$0.85 to
$\geq$0.99 on Bitcoin-Alpha and Bitcoin-OTC under a
\emph{transductive} evaluation regime.

We show that this trajectory does not reflect a corresponding
improvement in sign-prediction capability. The standard transductive
protocol leaks test-edge sign information into the model's feature
pipeline at training time, allowing every published method we audited
to recover the sign of held-out edges from features that should not
have seen it. The leakage is not a bug in any specific implementation
--- it is a structural property of how signed-graph message-passing
is currently set up. We diagnose it with a simple
\emph{label-shuffle audit}, define a \emph{strict training-edge-only
protocol} that closes the leak, re-evaluate the audited methods
under the corrected protocol, and present a compact architecture
--- Gömb-strict --- that remains competitive under the corrected
protocol with $30\,000$--$100\,000$ parameters and a total compute
budget of $\sim$24 GPU-hours on a single consumer GPU.

The corrected leaderboards differ from the published ones by
$[\text{Phase B}: \Delta_{\min}\text{--}\Delta_{\max}]$ percentage points
on average across methods and datasets; we recommend the
2018--2025 signed-link-prediction literature be re-evaluated under
the strict protocol.

% =========================================================
% § 2 — The mechanism (~400 words)
% =========================================================
\section{Why the transductive protocol leaks}

\label{sec:mechanism}

Signed graph neural networks construct per-vertex features by
aggregating signed messages from incident edges. In every method we
audited, the signed-incidence matrix
$\mathbf{H}_{\mathrm{signed}} \in \{-1, 0, +1\}^{|V|\times|E|}$ used
in this aggregation is constructed once from the \emph{full} edge
list, before the train/test split is applied. Concretely:

\begin{equation}
\mathbf{H}_{\mathrm{signed}}[u, e] = \begin{cases}
+1 & \text{if } u \in e \text{ and } \mathrm{sign}(e) = +1, \\
-1 & \text{if } u \in e \text{ and } \mathrm{sign}(e) = -1, \\
0  & \text{otherwise.}
\end{cases}
\end{equation}

The train/test split is then applied only to the \emph{supervisory
signal} --- which edges' signs are masked from the loss --- not to
the \emph{feature construction}. As a result, when a message-passing
layer computes vertex features as
$\mathbf{h}_v = \sigma\!\bigl(\sum_{e\ni v} \mathbf{H}_{\mathrm{signed}}[v, e]\,
\mathbf{W}\,\mathbf{h}_v^{(e)}\bigr)$,
the sign of every held-out test edge incident to $v$ is read directly
from $\mathbf{H}_{\mathrm{signed}}$ and propagated into $v$'s
representation. Backpropagation updates $\mathbf{W}$ using gradients
computed from these leaked features. At test time, the trained model
predicts the sign of a held-out edge $e^\star$ using vertex
representations that already contained $\mathrm{sign}(e^\star)$ as
an input.

The same mechanism applies to all derived signed-graph features:
signed Laplacians~\cite{sgcn2018}, balance-theoretic triad
indicators~\cite{cartwright1956}, and structural-equivalence
metrics. None of these can be computed without the full signed
adjacency, and none of the audited methods rebuilds them per training
mini-batch from training-edge signs only.

The fix is straightforward in principle and invasive in practice:
during training, every signed-feature artifact must be constructed
from training-edge signs alone, with test-edge contributions zeroed
out or replaced by a special "unknown" symbol. We call this the
\emph{strict training-edge-only protocol} and define it formally in
Section~\ref{sec:strict_protocol}.

A label-shuffle audit makes the leakage visible without committing
to any architectural fix: train under the standard protocol, then
randomly permute the test-edge signs (preserving the training-edge
signs), and measure test AUROC. If the feature pipeline is clean of
test-edge information, shuffled-test AUROC should equal chance
($0.5$ in expectation). If the feature pipeline leaks, shuffled-test
AUROC remains substantially above chance, in proportion to the
amount of test-edge sign information the features carry. The audit
is a falsifiable, model-agnostic test: any signed-link-prediction
method that fails it is, by definition, using information it
should not.

% ============================================================
% References (only what's cited above; coworker to merge with
% the manuscript's main .bib).
% ============================================================
\begin{thebibliography}{99}
  \bibitem{sgcn2018}    [TODO/CW] SGCN — Derr et al., ICDM 2018.
  \bibitem{sigat2019}   [TODO/CW] SiGAT — Huang et al., CIKM 2019.
  \bibitem{sgcl2022}    [TODO/CW] SGCL — Shu et al., AAAI 2022.
  \bibitem{sigformer2023} [TODO/CW] SiGformer — author/venue/year.
  \bibitem{sesgformer2024} [TODO/CW] SE-SGformer — author/venue/year.
  \bibitem{dadsgnn2025} [TODO/CW] DADSGNN — author/venue/year.
  \bibitem{cartwright1956} F. Cartwright \& D. Harary, "Structural
    balance: A generalization of Heider's theory,"
    \emph{Psychological Review} 63(5):277--293, 1956.
\end{thebibliography}

\end{document}
